představují pro \acacc{ČR} nejbližší negativní kontrolu. Sdílí postkomunistické dědictví, nízkou odborovou organizovanost, převahu podnikového vyjednávání a~obtížně využívané mechanismy rozšíření odvětvových smluv. I~zde zkušenost ukazuje, že samotná existence právního rámce nestačí -- pokud se odvětvové vyjednávání nestane prakticky použitelným nástrojem, pokrytí \aclgen{KS} zůstává nízké a~konvergence se snaží opřít hlavně o~nákladovou konkurenceschopnost.\par

V~komparativní perspektivě jde o~variantu, kde decentralizace \aclgen{KV} bez silné sektorové opory spíše prohlubuje roztříštěnost smluvního standardu. Data \aclgen{OECD} ukazují, že vysoké pokrytí přetrvává hlavně v~zemích se sektorovým nebo vícezaměstnavatelským vyjednáváním.~\cite{OECD_UnionEmployerCoverage_2025}~\cite{Hala_SystemySocialnihoDialogu2013}\par

Polsko bývá uváděno jako nejvýraznější růstový příběh Evropy posledních dvou dekád, jeho dynamika však nestojí na~kolektivním vyjednávání. Klíčovou roli sehrálo rychlé vybudování dopravní a~průmyslové infrastruktury -- spolufinancované evropskými kohezními fondy -- a~kvalifikovaná pracovní síla, které z~Polska učinily integrovaného subdodavatele německého průmyslu. Spojení blízkosti k~německému trhu, jednotného vnitřního trhu \aclgen{EU} a~vzdělané pracovní síly posunulo polskou ekonomiku od~levné práce k~průmyslové kapacitě s~vyšší přidanou hodnotou. Polsko bylo v~české konkurenci jako subdodavatel úspěšnější.\par

Na~tomto základě se~těžiště růstu přesunulo k~vnitřní poptávce. Podle hospodářské prognózy Evropské komise vzrostl reálný \acs{HDP} v~roce 2025 o~\SI{3,6}{\percent} tažen především soukromou spotřebou, přičemž růst nominálních náhrad zaměstnancům dosahoval kolem \SI{8}{\percent} při napjatém trhu práce s~nízkou nezaměstnaností kolem \SI{3}{\percent}. Polský případ tak ukazuje, že udržitelný růst může být tažen rostoucími mzdami a~na~ně navázanou agregátní spotřebou, v~Polsku se v~tento mzdový růst proměnilo napětí na trhu práce podobné \acs{ČR}.~\cite{EC_EconomicForecast_Poland_2026}\par

Slovensko se~od~\acgen{ČR} liší silnějším formálním ukotvením tripartity -- sociální dialog je uzákoněn prostřednictvím Rady hospodářské a~sociální dohody Slovenské republiky. Uzavírané generální dohody nemají právní, nýbrž pouze politickou a~morální závaznost, takže její plošný účinek zůstává omezený.\par

Od~roku 2011 může i~v~podnicích bez odborů uzavřít dohodu s~obsahem blízkým kolektivní smlouvě zaměstnanecká rada nebo zaměstnanecký důvěrník. Pokrytí kolektivními smlouvami se přesto na~Slovensku pohybuje přibližně kolem třetiny zaměstnanců, podobně jako v~\acloc{ČR}.~\cite{Hala_SystemySocialnihoDialogu2013}\par

Pro \acacc{ČR} z~toho plyne, že makroekonomická konvergence a~síla sociálního dialogu nejsou totéž, avšak kanál mzda--spotřeba--poptávka je pro udržitelný rozvoj zásadní. Zatímco Polsko jej naplnilo silnými subdodavatelskými vztahy, \ac{ČR} může tentýž efekt posílit institucionálně, rozšířením kolektivního vyjednávání jako stabilnějšího a~předvídatelnějšího mechanismu přenosu produktivity do~mezd.\par
